\setcounter{section}{1}
% Original source lines 69--78
\begin{Definition}
	Given a point $p$ in a Riemannian manifold $M$, and a collection $v_0,\dots,v_k$ of orthonormal vectors in $T_pM$, the $k^{\rm th}$-intermediate Ricci curvature at $p$ corresponding to this choice of vectors is defined to be $\sum_{i=1}^k K(v_0,v_i)$, where $K$ denotes the sectional curvature.
\end{Definition}
Notice that for an $n$-dimensional manifold, these curvatures interpolate between the Ricci curvature and the sectional curvature as $k$ ranges from $n-1$ down to 1. We will denote the $k^{\rm th}$-intermediate Ricci curvature by ${\rm Ric}_k$.

Our first main result is the following.

\begin{TheoremAlph}\label{main}
	Let $(M,g)$ be a compact Riemannian manifold, and suppose that $(M,g)\xrightarrow{\pi} (B^q,\check{g})$ is a Riemannian submersion with totally geodesic fibres. Suppose that $(B,\check{g})$ has ${\rm Ric}_{k_1}>0$, and that the fibres $\bigl(F_b^p,\hat{g}_b\bigr)=\bigl(\pi^{-1}(b),g|_{\pi^{-1}(b)}\bigr)$ for $b \in B$ $($which are necessarily isometric$)$ have ${\rm Ric}_{k_2}>0$. Then there exists $\tau=\tau(M,g,\pi)>0$ such that for all positive $t<\tau$, scaling $g$ in fibre directions by $t$ yields a metric on $M$ with ${\rm Ric}_k>0$ for every $k\ge \max\{k_1+p,k_2+q\}$.	
\end{TheoremAlph}


% Original source lines 128--138
In order to establish Theorem \ref{main} and Corollary \ref{fibre_bundles}, we study the following operator. For $X \in T_pM$, define
\begin{align*}
	R_X\colon \ T_pM &\to T_pM, \\
	Y &\mapsto R(X,Y)X
\end{align*}
(Note that, depending on the chosen sign convention for the curvature tensor, there exist different conventions to define this operator. See Section \ref{formulas} for our choice of sign convention.)
It is easy to see that this operator is symmetric, so there exists an orthonormal basis of eigenvectors. Given orthonormal unit vectors $X$, $Y$, if $Y$ is an eigenvector of $R_X$ with eigenvalue $\lambda$, then \[\lambda=\langle R_X(Y),Y\rangle=K(X,Y).\] Thus we can obtain $k^{\rm th}$ intermediate Ricci curvatures by summing eigenvalues of $R_X$. More precisely, we have the following.
\begin{Lemma}\label{k+1}
	A Riemannian manifold $M$ has ${\rm Ric}_k>0$ if and only if for every unit vector $X \in TM$, the sum of any $(k+1)$-eigenvalues of $R_X$ is positive.
\end{Lemma}
Note that this is essentially \cite[Lemma 2.6]{RW}, and we refer the reader there for more details.~(Of course $X$ is an eigenvector of $R_X$ with eigenvalue~0, and therefore $\langle R_X(X),X\rangle =K(X,X)=0$, hence the need to sum $k+1$ eigenvectors to obtain a positive sum.) Theorem \ref{main} is then obtained by a careful analysis of the eigenvectors of $R_X$.


% Original source lines 142--289
\section{Curvature formulae}\label{formulas}

In this section, we present some formulas for the curvature tensor of a Riemannian submersion. We are interested in the case where the fibres are totally geodesic, and the effect the canonical variation has on the curvature tensor in this case. Recall that the canonical variation involves scaling the fibre metrics by a uniform constant, leaving the base metric and the horizontal distribution untouched. The case of particular interest is when the fibre scaling factor approaches zero. Note that the canonical variation preserves the property of being totally geodesic.

The starting point for the curvature formulas we will require is the O'Neill formulas for Riemannian submersions, as presented, for example, in~\cite{Be}. As this book is the definitive reference for the curvature of Riemannian submersions and the canonical variation, we will adopt the conventions used there. In particular, note that for the canonical variation, we will scale in fibre directions by a constant $t>0$ \big(as opposed to the more usual $t^2$\big). We also use the following curvature tensor convention:
\[
\langle R(V,W)X,Y\rangle=\bigl\langle \nabla_W\nabla_V X -\nabla_V \nabla_W X+\nabla_{[V,W]}X,Y\bigr\rangle,
\]
where $\langle\,,\,\rangle$ denotes the Riemannian metric.

The formulas below are due to O'Neill and Gray. They appear in \cite{Be} as Theorem 9.28. In~stating these, we have specialized to the case of totally geodesic fibres, i.e., we have assumed that the $T$-tensor is identically zero. Recall that the \emph{$A$-tensor} of a Riemannian submersion is a~$(2,1)$-tensor field which vanishes if and only if the horizontal distribution is integrable. We refer to \cite[Section 9.C]{Be} for its definition and basic properties.

\begin{Lemma}\label{formulas1}
	Given a Riemannian submersion with totally geodesic fibres, let $X$, $Y$, $Z$, $Z'$ denote arbitrary horizontal vectors, and $U$, $V$, $W$, $W'$ arbitrary vertical vectors, all over the same base point. Then
	\begin{align*}
		&\bigl\langle R(U,V)W,W' \bigr\rangle =\bigl\langle \hat{R}(U,V)W,W' \bigr\rangle, \\
		&\langle R(U,V)W,X \rangle =0, \\
		&\langle R(X,U)Y,V \rangle =\langle (\nabla_U A)_X Y,V \rangle+\langle A_X U, A_Y V\rangle, \\
		&\langle R(U,V)X,Y \rangle =\langle (\nabla_U A)_X Y,V \rangle-\langle (\nabla_V A)_X Y,U\rangle+\langle A_XU,A_YV \rangle-\langle A_XV, A_YU \rangle, \\
		&\langle R(X,Y)Z,U \rangle = \langle (\nabla_ZA)_X Y,U \rangle, \\
		&\bigl\langle R(X,Y)Z,Z' \bigr\rangle =\bigl\langle \check{R}(X,Y)Z,Z' \bigr\rangle-2\bigl\langle A_XY,A_ZZ' \bigr\rangle + \bigl\langle A_Y Z,A_X Z'\bigr\rangle-\bigl\langle A_X Z, A_Y Z'\bigr\rangle.
	\end{align*}
	Here, $\hat{R}$ is the curvature tensor of the fibre metric, and $\check{R}$ that of the base metric.
\end{Lemma}

The next lemma is an adaptation of \cite[Lemma 9.69]{Be}, where we only state the effect of the canonical variation on formulas which appear in Lemma \ref{formulas1}.

\begin{Lemma}\label{formulas2}
	Given a Riemannian submersion with fibre metric $\hat{g}$, we scale in fibre directions by $t>0$, that is, we take the new fibre metric to be $t\hat{g}$, but leave the base metric and horizontal distribution unchanged. Denoting quantities for the rescaled metric with a superscript $t$, and using the same vector notation as in Lemma~$\ref{formulas1}$, the following formulas hold:
	\begin{align*}
		&A^t_X Y=A_XY, \\
		&A^t_X U=tA_XU, \\
		&[\langle (\nabla_X A)_YZ,U\rangle]^t=t\langle (\nabla_XA)_Y Z,U\rangle, \\
		&[\langle (\nabla_UA)_X Y,V\rangle]^t=t\langle (\nabla_U A)_X Y,V \rangle+\bigl(t-t^2\bigr)\bigl(\langle A_XU,A_YV\rangle-\langle A_XV,A_YU\rangle\bigr).
	\end{align*}
\end{Lemma}

It can be deduced from the Koszul formula, for example, that scaling a metric $g$ by a constant factor, i.e., replacing $g$ by $t\cdot g$ for $t>0$, leaves the Levi-Civita connection unchanged. Therefore given vectors $U$, $V$, $W$, we see immediately that $R(U,V)W$ is unaffected by scaling the metric as well. Hence, in the case of arbitrary vertical vectors $U$, $V$, $W$, $W'$, we then have
\[
\bigl[\bigl\langle\hat{R}(U,V)W,W'\bigr\rangle\bigr]^t=t\bigl\langle\hat{R}(U,V)W,W'\bigr\rangle.
\]
Further, since the metric is unchanged in horizontal directions, for horizontal vectors~$X$,~$Y$,~$Z$,~$Z'$, we have
\[
\bigl[\bigl\langle \check{R}(X,Y)Z,Z'\bigr\rangle\bigr]^t=\bigl\langle \check{R}(X,Y)Z,Z'\bigr\rangle.
\]
Combining these observations with Lemma \ref{formulas2}, a computation yields the following formulas, which show how the expressions Lemma \ref{formulas1} transform under the canonical variation.

\begin{Corollary}\label{formulas3}
	Under the same assumptions and with the same notation as the previous two lemmas, the following formulas hold:
	\begin{align*}
		&\bigl[\bigl\langle R(U,V)W,W' \bigr\rangle\bigr]^t =t\bigl\langle \hat{R}(U,V)W,W' \bigr\rangle, \\
		&[\langle R(U,V)W,X\rangle]^t =0, \\
		&[\langle R(X,U)Y,V \rangle]^t =t\langle (\nabla_U A)_X Y,V \rangle+\bigl(t-t^2\bigr)(\langle A_XU,A_YV\rangle-\langle A_XV,A_YU\rangle)\\
		&\hphantom{[\langle R(X,U)Y,V \rangle]^t =}{} +t^2\langle A_X U, A_Y V\rangle \\
		&\hphantom{[\langle R(X,U)Y,V \rangle]^t }{} =t\langle (\nabla_U A)_X Y,V \rangle +t\langle A_X U, A_Y V\rangle+\bigl(t^2-t\bigr)\langle A_XV,A_YU\rangle,\\
		&[\langle R(U,V)X,Y \rangle]^t =t\langle (\nabla_U A)_X Y,V \rangle-t\langle (\nabla_V A)_X Y,U\rangle+t^2\langle A_XU,A_YV \rangle-t^2\langle A_XV, A_YU \rangle \\
		&\hphantom{[\langle R(U,V)X,Y \rangle]^t =}{} +\bigl(t-t^2\bigr)(\langle A_XU,A_YV\rangle-\langle A_XV,A_YU\rangle\\
&\hphantom{[\langle R(U,V)X,Y \rangle]^t =}{}- \langle A_XV,A_YU\rangle + \langle A_XU,A_YV\rangle)\\
		&\hphantom{[\langle R(U,V)X,Y \rangle]^t }{} =t\langle (\nabla_U A)_X Y,V \rangle-t\langle (\nabla_V A)_X Y,U\rangle\\
		&\hphantom{[\langle R(U,V)X,Y \rangle]^t =}{} +\bigl(2t-t^2\bigr)(\langle A_XU,A_YV\rangle-\langle A_XV,A_YU\rangle),\\
		&[\langle R(X,Y)Z,U \rangle]^t = t\langle (\nabla_ZA)_X Y,U \rangle, \\
		&\bigl[\bigl\langle R(X,Y)Z,Z' \bigr\rangle\bigr]^t=\bigl\langle \check{R}(X,Y)Z,Z' \bigr\rangle\!-2t\bigl\langle A_XY,A_ZZ' \bigr\rangle\! + t\bigl\langle A_Y Z,A_X Z'\bigr\rangle\!-t\bigl\langle A_X Z, A_Y Z'\bigr\rangle.
	\end{align*}
\end{Corollary}

Let us fix the following notation. Let $g$ denote the original (unscaled) Riemannian submersion metric, and denote by $g_t$ the canonical variation metric corresponding to the vertical scaling by~$t$. We now wish to specialize to the case where all vectors are unit vectors for $g_t$. Suppose that~$U$,~$V$,~$W$,~$W'$ are all unit vertical vectors for $g$ (i.e., the case $t=1$), then the corresponding unit vectors for $g_t$ are $U/\sqrt t$, $V/\sqrt t$, $W/\sqrt t$, $W'/\sqrt t$. From Corollary \ref{formulas3}, we immediately obtain the expressions below.

\begin{Corollary}\label{formulas4}
	Suppose that $X$, $Y$, $Z$, $Z'$ are unit horizontal vectors for $g_t$. If $U$, $V$, $W$, $W'$ are all unit vertical vectors for $g$, and $U_t$, $V_t$, $W_t$, $W'_t$ the corresponding unit vectors for $g_t$, we~have
	\begin{align*}
		&\bigl[\bigl\langle R(U_t,V_t)W_t,W'_t \bigr\rangle\bigr]^t =\frac{1}{t}\bigl\langle \hat{R}(U,V)W,W' \bigr\rangle, \\
		&[\langle R(U_t,V_t)W_t,X\rangle]^t =0, \\
		&[\langle R(X,U_t)Y,V_t \rangle]^t =\langle (\nabla_U A)_X Y,V \rangle +\langle A_X U, A_Y V\rangle+(t-1)\langle A_XV,A_YU\rangle,\\
		&[\langle R(U_t,V_t)X,Y \rangle]^t =\langle (\nabla_U A)_X Y,V \rangle-\langle (\nabla_V A)_X Y,U\rangle\\
		&\hphantom{[\langle R(U_t,V_t)X,Y \rangle]^t =}{} +(2-t)(\langle A_XU,A_YV\rangle-\langle A_XV,A_YU\rangle),\\
		&[\langle R(X,Y)Z,U_t \rangle]^t = \sqrt{t}\langle (\nabla_ZA)_X Y,U \rangle, \\
		&\bigl[\bigl\langle R(X,Y)Z,Z' \bigr\rangle\bigr]^t =\bigl\langle \check{R}(X,Y)Z,Z' \bigr\rangle\!-2t\bigl\langle A_XY,A_ZZ' \bigr\rangle\! + t\bigl\langle A_Y Z,A_X Z'\bigr\rangle\!-t\bigl\langle A_X Z, A_Y Z'\bigr\rangle.
	\end{align*}
\end{Corollary}

\begin{Corollary}\label{R_matrix}
	Let $\Psi$ be a unit tangent vector with respect to $g_t$, i.e., we can write $\Psi=\lambda U_t+\mu X$ with $\lambda,\mu\in\mathbb{R}$ so that $\lambda^2+\mu^2=1$. The map $R_\Psi^t$, according to the splitting into vertical and horizontal parts, is given by the matrix
	\begin{equation*}
	\begin{pmatrix}
		\frac{\lambda^2}{t}\hat{R}_U+O(t) & O(|\lambda|)+O\bigl(\sqrt{t}\bigr)\\
		O(|\lambda|)+O\bigl(\sqrt{t}\bigr) & \mu^2 \check{R}_X+O(|\lambda|)+O(t)
	\end{pmatrix}
\end{equation*}
	as $t,\lambda\to0$.
\end{Corollary}
Here for a linear map $T_\varepsilon$ that depends on a parameter $\varepsilon$ we say that $T_\varepsilon=O(f(\varepsilon))$ for a real function $f$, if the operator norm $\lVert T_\varepsilon\rVert$ satisfies $\lVert T_\varepsilon\rVert=O(f(\varepsilon))$ in the usual sense.
\begin{proof}
	We first consider the first entry of the matrix. Here, we have
	\begin{align*}
		[\langle R(\Psi,V_t)\Psi,W_t\rangle]^t&{}=\lambda^2[\langle R(U_t,V_t)U_t,W_t\rangle]^t+\mu^2 [\langle R(X,V_t)X,W_t\rangle]^t\\
		&\quad{} + \lambda\mu\bigl([\langle R(U_t,V_t)X,W_t\rangle]^t + [\langle R(X,V_t)U_t,W_t\rangle]^t \bigr).
	\end{align*}
	By Corollary \ref{formulas4}, we have
	\begin{align*}
		[\langle R(U_t,V_t)U_t,W_t\rangle]^t &{}= \frac{1}{t}\bigl\langle \hat{R}(U,V)U,W\bigr\rangle,\\
		[\langle R(X,V_t)X,W_t\rangle]^t &{}= \langle (\nabla_V A)_X X,W\rangle +t\langle A_X V,A_X W\rangle,\\
		[\langle R(U_t,V_t)X,W_t\rangle]^t &{}=0,\\
		[\langle R(X,V_t)U_t,W_t\rangle]^t &{}=0.
	\end{align*}
	It can be checked that $\nabla_V A$ enjoys the same antisymmetry property with respect to pairs of horizontal vectors as the $A$-tensor itself. Thus $\langle (\nabla_V A)_X X,W\rangle=0$. Since the $A$-tensor is bounded (as is any operator on a finite-dimensional vector space) and since $|\mu|\leq 1$, the first entry of the matrix follows.
	
	Similarly, for the second entry we have
	\begin{align*}
		[\langle R(\Psi,V_t)\Psi,Y\rangle ]^t&{}= \lambda^2[\langle R(U_t,V_t)U_t,Y\rangle]^t+\mu^2 [\langle R(X,V_t)X,Y\rangle]^t\\
		&\quad{} + \lambda\mu\bigl([\langle R(U_t,V_t)X,Y\rangle]^t + [\langle R(X,V_t)U_t,Y\rangle]^t \bigr).
	\end{align*}
	Again, by Corollary \ref{formulas4}, we have
	\begin{align*}
		[\langle R(U_t,V_t)U_t,Y\rangle]^t&{}=0,\\
		[\langle R(X,V_t)X,Y\rangle]^t&{}=\sqrt{t}\langle(\nabla_X A)_X Y,V\rangle=O\bigl(\sqrt{t}\bigr),
	\end{align*}
	and both $[\langle R(U_t,V_t)X,Y\rangle]^t$ and $[\langle R(X,V_t)U_t,Y\rangle]^t$ are bounded as $t\to 0$. Hence, we obtain, again by using $|\mu|\leq 1$,
	\begin{align*}
		[\langle R(U_t,V_t)U_t,Y\rangle]^t&{}=O\bigl(\sqrt{t}\bigr)+O(|\lambda|).
	\end{align*}
	By the symmetry properties of $R_\Psi$, we obtain the same conclusion for the third entry.
	
	Finally, for the fourth entry, we have
	%\begin{align*}
	%		[\langle R(\Psi,Y)\Psi,V_t\rangle ]^t=& [\langle R(\Psi,V_t)\Psi,Y\rangle ]^t=O(|\lambda|)+O\bigl(\sqrt{t}\bigr),\\ \\
	%\end{align*}
	\begin{align*}
		[\langle R(\Psi,Y)\Psi,Z\rangle ]^t&{}= \lambda^2[\langle R(U_t,Y)U_t,Z\rangle]^t+\mu^2 [\langle R(X,Y)X,Z\rangle]^t\\
		&\quad{}+ \lambda\mu\bigl([\langle R(U_t,Y)X,Z\rangle]^t + [\langle R(X,Y)U_t,Z\rangle]^t \bigr)
		%		&{}= O(|\lambda|)+\mu^2(\langle \check{R}(X,Y)X,Z\rangle+O(t) )\\
		%		&{}= \mu^2(\langle \check{R}(X,Y)X,Z\rangle)+O(|\lambda|)+O(t).
	\end{align*}		
	By Corollary \ref{formulas4}, the terms $[\langle R(U_t,Y)U_t,Z\rangle]^t$, $[\langle R(U_t,Y)X,Z\rangle]^t$, and $[\langle R(X,Y)U_t,Z\rangle]^t$ are bounded as $t\to 0$ and
	\[ [\langle R(X,Y)X,Z\rangle]^t=\bigl\langle\check{R}(X,Y)X,Z\bigr\rangle+O(t). \]
	Hence, using $|\mu|,|\lambda|\leq 1$, we obtain
	\begin{align*}
		[\langle R(\Psi,Y)\Psi,Z\rangle ]^t&{}= O(|\lambda|)+\mu^2\bigl(\bigl\langle \check{R}(X,Y)X,Z\bigr\rangle+O(t) \bigr)\\
		&{}= \mu^2\bigl(\bigl\langle \check{R}(X,Y)X,Z\bigr\rangle\bigr)+O(|\lambda|)+O(t).
 \tag*{\qed}
 \end{align*}
 \renewcommand{\qed}{}
\end{proof}



\section{Proof of the main theorems}\label{proofs}




% Original source lines 292--323
\begin{proof}[Proof of Theorem \ref{main}]
	We fix a point $p\in M$. As in Corollary \ref{R_matrix}, any unit tangent vector $\Psi\in T_p M$ with respect to $g_t$ can be written as $\Psi=\lambda U_t+\mu X$ with $\lambda^2+\mu^2=1$ and $\lambda,\mu\geq 0$. In~the~following we therefore fix $U$ and $X$, consider $\Psi=\lambda U_t+\mu X$ as parametrized by $\lambda $ and $t$, and show that for all $t$ sufficiently small and all $\lambda\in[0,1]$, the sum of any $(k+1)$ eigenvalues of the map $R^t_\Psi$ is positive for $k\geq\max\{k_1+p,k_2+q \}$. (Recall that by Lemma~\ref{k+1}, this is equivalent to the condition ${\rm Ric}_k>0$.)
	
	By Corollary \ref{R_matrix}, the map $R_\Psi^t$ is given by the matrix
	\[
	\begin{pmatrix}
		\frac{\lambda^2}{t}\hat{R}_U+O(t) & O(|\lambda|)+O\bigl(\sqrt{t}\bigr)\\
		O(|\lambda|)+O\bigl(\sqrt{t}\bigr) & \mu^2 \check{R}_X+O(|\lambda|)+O(t)
	\end{pmatrix}
	\]
	as $t,\lambda\to0$. For $\lambda=0$, we obtain the matrix
	\[
	\begin{pmatrix}
		O(t) & O\bigl(\sqrt{t}\bigr)\\
		O\bigl(\sqrt{t}\bigr) & \check{R}_X+O(t)
	\end{pmatrix}.
	\]
	%Recall that by Lemma \ref{k+1}, we have ${\rm Ric}_k>0$ on $M$ if and only if for each unit $\Psi$ the map $R_\Psi^t$ is $(k+1)$-positive.
	By assumption $\check{R}_X$ is $k_1$-positive, so the sum of any $k_1+1$ of its eigenvalues is positive. Since eigenvalues depend continuously on the matrix, see, e.g., \cite[Corollary VI.1.6]{Bh}, there exists $t_0>0$ so that the sum of any $p+k_1+1$ eigenvalues of the above matrix is positive for all $t\in(0,t_0)$. Since $R^t_\Psi$ depends continuously on $\Psi$, this conclusion holds for all $\lambda\in[0,\lambda_0)$ for some $\lambda_0>0$.
	
	Now let $\lambda\geq\lambda_0$. Multiplying the matrix by $t$ gives the matrix
	\[
	\begin{pmatrix}
		\lambda^2 \hat{R}_U+O\bigl(t^2\bigr) & t\cdot O(|\lambda|)+O\bigl(t^{\frac{3}{2} }\bigr)\\
		t\cdot O(|\lambda|)+O\bigl(t^{\frac{3}{2}}\bigr) & \mu^2t \check{R}_X+t\cdot O(|\lambda|)+O\bigl(t^2\bigr)
	\end{pmatrix},
	\]
	and note that the condition of $k$-positivity is invariant under scalar multiplication of the matrix by a positive constant. As the fibre is compact, $\hat{R}_U$ is $(k_2+1)$-bounded away from zero. Thus since $\lambda\in[\lambda_0,1]$, and the $O(|\lambda|)$-terms depend continuously on $\lambda$ and hence are bounded for $\lambda\in[\lambda_0,1]$, there exists $t_1>0$ so that the sum of any $k_2+q+1$ eigenvalues of this matrix is positive for all $t\in(0,t_1)$. %Again by continuity, $t_1$ can be chosen to be independent of $U,X$ and $\mu$.
	
	We set $\tau=\min(t_0,t_1)$, so the sum of any $k+1$ eigenvalues of $R_\Psi^t$ is positive for all $t\in(0,\tau)$ and $\lambda\in[0,1]$, and for all $k\geq \max\{k_1+p,k_2+q\}$. By continuity and compactness of the space of unit horizontal respectively vertical vectors, it then follows that there exists a uniform choice of $t$ for all unit tangent vectors $\Psi$, and by compactness of $M$, for all $p\in M$.
	%for all unit tangent vectors $\Psi\in T_p M$ \green{and all $k\geq \max\{k_1+p,k_2+q\}$}. By openness of the ${\rm Ric}_k>0$ condition, this conclusion holds in a neighbourhood of $p$. Hence, since $B$ is compact and $b$ is chosen arbitrary, the conclusion holds globally after possibly further decreasing $\tau$.
\end{proof}
